% Source-derived chapter 08: Proofs of the main theorems
% Original source: riemann-roch-polynomials-hyperkahler-fourfolds
\chapter{Proofs of the main theorems}
\label{riemann-roch-polynomials-hyperkahler-fourfolds-chapter-08}
\source{riemann-roch-polynomials-hyperkahler-fourfolds}

\begin{remark}[Source section]
\label{riemann-roch-polynomials-hyperkahler-fourfolds-chapter-08-source}
\source{riemann-roch-polynomials-hyperkahler-fourfolds}
\source{riemann-roch-polynomials-hyperkahler-fourfolds}
This chapter reproduces the corresponding section of the retrieved
LaTeX source, including its definitions, statements, examples, and
proofs. It has not yet been translated into Lean.
\notready
\end{remark}

% The source section title was: Proofs of the main theorems
\label{sec:OGrady}
\source{riemann-roch-polynomials-hyperkahler-fourfolds}

In this section, we prove Theorem~\ref{thm:MainThm} and Theorem~\ref{thm:SYZ}.\
Let $X$ be a  \hk\ fourfold.
We assume that there are classes $\lll,\mm\in H^2(X,\Z)$ with $\int_X \lll^4=0$ and $\int_X\lll^2\mm^2=2$.

The sublattice spanned by $\lll$ and $\mm$ is indefinite, hence its $q_X$-orthogonal has signature $(2,b_2(X)-4)$.\ 
By the surjectivity of the period map (\cite{HuyHK}), we can then deform $X$ and assume that  $(X,\lll,\mm)$ is a very general triple satisfying~\eqref{hypo} (hence also~(\ref{hypo}$^\prime$)).\
After possibly permuting~$\lll$ and $\mm$, we can further assume that $\lll$ is
nef (see~\eqref{cones}).\
By Corollary~\ref{cor:NoTwoIsotropicNef}, the class $\mm$ cannot be nef, and we are in the situation of Proposition~\ref{prop:CaseC2}.\
Let us denote by $f\colon X \to \P^2$ the Lagrangian fibration such that $f^*\cO_{\P^2}(1)\isom L^{\otp k_L}$, for some positive integer $k_L$.\
We start with the following vanishing result.

\begin{lemm}\label{lem:vanishing}
\source{riemann-roch-polynomials-hyperkahler-fourfolds}
Under the above assumptions, let $p,q\in\Z$, with $q>0$.\
For all $i\geq 3$, we have
\[
H^i(X,L^{\otp p}\otimes M^{\otp q})=0.
\]
\end{lemm}

\begin{proof}
Choose $r\in\Z$ be such that $p+rk_L \geq q$.\
We can write
\[
L^{\otp p}\otimes M^{\otp q} =  L^{\otp p+rk_L}  \otimes L^{-rk_L} \ot M^{\otp q}=  L^{\otp p+rk_L} \otimes M^{\otp q}\otimes f^*\cO_{\P^2}(-r),
\]
where $ L^{\otp p+rk_L}\ot M^{\otp q}$ is big and nef.\
By \cite[Theorem~10.32]{kol} again, we obtain that $R^j f_* (L^{\otp p}\otimes M^{\otp q})$ vanishes for all $j>0$ and $f_* (L^{\otp p}\otimes M^{\otp q})$ is a vector bundle on~$\P^2$.\ 
In particular, this gives us what we need:
\[
H^i(X, L^{\otp p}\otimes M^{\otp q}) = H^i(\P^2, f_* (L^{\otp p}\otimes M^{\otp q})) = 0,
\]
for all $i\geq3$.
\end{proof}

We now study in more detail the surface $\Sigma$ and the vector bundle $\cE$.\
Recall that we have $(L\otimes M)\vert_E = c_E^* H_{\Sigma}$.\
Let us denote by $\hh\in \NS(\Sigma)$ the  class of $H_{\Sigma}$.

\begin{lemm}\label{lem:K3}
\source{riemann-roch-polynomials-hyperkahler-fourfolds}
The pair $(\Sigma,\hh)$ is a polarized K3 surface of degree~2 with $\NS(\Sigma)=\Z \hh$.
\end{lemm}

\begin{proof}
By~\cite[Theorem~1.4]{wie}, we know that $\Sigma$ is a symplectic surface.\
To show that it is a K3 surface, it is enough to compute $\chi(\Sigma,\cO_\Sigma)=\chi(E,\cO_E)$.\
By the Riemann--Roch Theorem, we have
\[
\chi(E,\cO_E) = \chi(X,\cO_X)-\chi(X,\cO_X(-E))=3-1=2,
\]
as we wanted.\
Also,
\[
\hh^2 = \int_X (\lll+\mm)^2 (-\lll+\mm) \lll = 2.
\]
Thus, we are left to show that the N\'eron--Severi group of $\Sigma$ has rank one.



Let us consider the transcendental lattice $H^2(X,\Z)_{\rm tr}\coloneqq \NS(X)^\perp\subset H^2(X,\Z)$, which under our assumptions has rank~21.
For all $\alpha\in H^2(X,\Z)_{\rm tr}$,  we have, by~\eqref{eq:Fujiki3},
\[
\int_X \alpha(-\lll+\mm)(\lll+\mm)^2 = 0.
\]
Hence, we can write
\[
\alpha\vert_E = c_E^*\nn_{\alpha}
\]
for a unique class $\nn_{\alpha}\in \hh^\bot \subset H^2(\Sigma,\Z)$.
Moreover, again by~\eqref{eq:Fujiki3}, for all $\alpha,\beta\in H^2(X,\Z)_{\rm tr}$, we have
\[
q_X(\alpha,\beta) = \int_X \alpha\beta\lll(-\lll+\mm) = \int_{\Sigma}\nn_{\alpha}\nn_{\beta}.
\]
Therefore, the morphism
\[
\vartheta\colon H^2(X,\Z)_{\rm tr} \lra \hh^\bot \subset H^2(\Sigma,\Z), \qquad \alpha\longmapsto \nn_{\alpha}
\]
gives an isometry $\vartheta_{\C}\colon H^2(X,\Z)_{\rm tr}\otimes\C \isomto \hh^\perp\otimes\C$.

The morphism $\vartheta$ is a nonzero morphism of Hodge structures.\ Since the Hodge structure~$H^2(X)_{\rm tr}$
is irreducible,~$\vartheta$ is injective  hence $\Sigma$ has Picard number one.
 \end{proof}

\begin{lemm}\label{lem:bundle}
\source{riemann-roch-polynomials-hyperkahler-fourfolds}
The vector bundle $\cE$ is the unique spherical stable bundle on $\Sigma$ with Mukai vector~$(2,H_\Sigma,1)$.
\end{lemm}

\begin{proof}
By Lemma~\ref{lem:K3}, the Picard group of $\Sigma$ is generated by $H_\Sigma$.\
Hence, we can write  the Mukai vector of $\cE$ as $v(\cE)=(2,s H_\Sigma, s')$, with $s,s'\in\Z$.\
By the Riemann--Roch Theorem, since $[E]=-\lll+\mm$, we have
\[
\chi(X,L (-E))=\chi(X,L^{\otp 2}\otimes M^{-1})=0.
\]
Hence, from the exact sequence
\begin{equation}\label{eq:Paris20211029}
\source{riemann-roch-polynomials-hyperkahler-fourfolds}
0\to L (-E) \to L \to L \vert_E \to 0
\end{equation}
we get
\[
s'+2 = \chi(\Sigma,\cE) = \chi(E,L \vert_E) = \chi(X,L) - \chi(X,L (-E)) = \chi(X,L) = 3,
\]
and thus $s'=1$.
To compute $s$, we proceed similarly and use the exact sequence
\begin{equation}\label{eq:Paris20211028}
\source{riemann-roch-polynomials-hyperkahler-fourfolds}
0\to M^{-1}  (-E) \to M^{-1} \to M^{-1} \vert_E \to 0.
\end{equation}
Again, since $\cO_X(-E)=L \otimes M^{-1}$, we get
\[
M^{-1} (-E) = L\otimes M^{-2}.
\]
Hence
\begin{equation}\label{eq:bundle}
\source{riemann-roch-polynomials-hyperkahler-fourfolds}
\begin{split}
5 - 2s & = \chi(\Sigma,\cE \otimes H_\Sigma^{-1}) = \chi(E,M^{-1} \vert_E) \\
& = \chi(X,M^{-1}) - \chi(X,L\otimes M^{-2}) = \chi(X,M^{-1}) = 3,
\end{split}
\end{equation}
and thus $s=1$.

To finish the proof, we only need to prove that $\cE$ is stable; indeed, it is then automatically spherical, since its Mukai vector has square  $-2$.
Since $\cE$ is a rank-2 vector bundle of slope $(H_\Sigma\cdot c_1(\cE))/(H_\Sigma^2)\rk(\cE)=1/2$ and $\Sigma$ has Picard number 1, it is enough to prove   \mbox{$H^0(\Sigma,\cE\otimes H_\Sigma^{-1})=0$.}\
This is a slight refinement of~\eqref{eq:bundle} above.\
Indeed,
\[
H^0(\Sigma,\cE\otimes H_\Sigma^{-1})=H^0(E,(L\otimes (L\otimes M)^{-1})\vert_E)=H^0(E,M^{-1} \vert_E).
\]
From~\eqref{eq:Paris20211028}, we obtain an exact sequence
\[
H^0(X,M^{-1}) \to H^0(E,M^{-1} \vert_E) \to H^1(X,L\otimes M^{-2}).
\]
Under our assumptions, $\mm$ is contained in the closure of the positive cone and therefore $-\mm$ cannot be effective; hence, $H^0(X,M^{-1})=0$.\
Therefore, we only have to show  $H^1(X,L\otimes M^{-2})=0$ or, by Serre duality,  $H^3(X,  L^{-1}\ot M^{\otp 2})=0$.
This follows immediately from Lemma~\ref{lem:vanishing}.
\end{proof}

Next, we use Lemma~\ref{lem:bundle} to show $k_L=1$ and study the induced map $f_E\coloneqq f \vert_E \colon E\to \P^2$.

\begin{lemm}\label{lem:k=1}
\source{riemann-roch-polynomials-hyperkahler-fourfolds}
The restriction morphism
\[
H^0(X,L)\lra H^0(E, L \vert_E) \isom \C^3
\]
is an isomorphism.
In particular, $k_L=1$ and $L$ is globally generated.
\end{lemm}

\begin{proof}
We consider again the exact sequence~\eqref{eq:Paris20211029}.\
By Lemma~\ref{lem:vanishing} and Serre duality, we get
\[
h^i(X,L (-E)) = h^i(X,L^{\otp 2}\otimes M^{-1})=h^{4-i}(X, L^{-2}\otimes M)=0,
\]
for $i\in\{0,1\}$, as we wanted.
The equality $h^0(E, L \vert_E)=3$ and the last statement follow since, by the Riemann--Roch Theorem, we have that $h^0(E, L \vert_E)\geq 3$, and then we apply Lemma~\ref{lem:H0large}.
\end{proof}

The moduli space $\M_0(\Sigma)$ was defined in Example~\ref{ex:K3n}.\
Let us denote by $E_0$ the exceptional divisor given by the universal family of curves in $\Sigma$ over $\P^2=|\hh |$.
Then $E_0$ is isomorphic to $ E$ over $\Sigma$ since they are projective bundles associated with  the same vector bundle $\cE$.

To finish the proof of Theorem~\ref{thm:MainThm} and Theorem~\ref{thm:SYZ}, we only need to show the following result, which can be seen as a version of~\cite[Theorem 1]{nag}.

\begin{prop}
\source{riemann-roch-polynomials-hyperkahler-fourfolds}
\notready\label{thm:OGrady}
\source{riemann-roch-polynomials-hyperkahler-fourfolds}
Let $(X,\lll,\mm)$ be a very general triple satisfying~\eqref{hypo} with $\mm$ isotropic and $\lll$ nef.\
Then $X$ is isomorphic to $\M_0(\Sigma)$.\
In particular, $X$ is of $\mathrm{K3}^{[2]}$ deformation type.
\end{prop}

\begin{proof}
By Lemma~\ref{lem:k=1}, the morphism $f_E$ coincides with the composition
\[
E \isomlra E_0 \lhra \M_0(\Sigma) \xrightarrow{\ f_0\ } \P^2,
\]
and thus $E$ is isomorphic to the universal family of curves in $\Sigma$ over $\P^2=|\hh |$.
Let us consider the nonempty open subset $U\subset\P^2$ where all morphisms $f$, $f_0$ and their restrictions respectively to $E$ and $E_0$ are smooth.
Then the restriction $f_E \vert_{U_E}$ of $f_E$ to $U_E\coloneqq f_E^{-1}(U)$ will factor via the relative Albanese variety $\Alb (U_E/U)$ (see~\cite[Proposition 1]{fujiki}): there exists a proper morphism
\[
g\colon \Alb (U_E/U) \lra f^{-1}(U)
\]
over $U$ compatible with the inclusion of $U_E$.

For all points $u\in U$, the class of the curve $C_u\coloneqq f_E^{-1}(u)$ gives a principal polarization on the abelian surface $A_u\coloneqq f^{-1}(u)$.
Since $C_u$ lives in a very general K3 surface, $A_u$ cannot be isomorphic to the product of two elliptic curves.\
Hence, $A_u$ is isomorphic to the Jacobian of~$C_u$ and the principal polarization is the theta polarization.\
From this, we immediately deduce that $g$ is an isomorphism.\
But the relative Albanese variety only depends on $f_E$ and thus it is isomorphic to the relative Albanese variety of the universal family of curves in $\Sigma$ over $\P^2=|\hh |$.\
This gives a birational morphism $X \isomdra  \M_0(\Sigma)$ which is then an isomorphism, since the nef and movable cones coincide.
\end{proof}


%%%%%%%%%%%%%%%%%%%%%
